% Source: https://artofproblemsolving.com/community/c3h378601_mathcounts_state_2004_4
% Stable identity: aops-topic-378601-q04
The other leg is $\sqrt{97-16}=\sqrt{81}=9$ inches, so the area is
\[
\dfrac{1}{2}\cdot9\cdot4=18\text{ square inches}.
\]
A square foot is $12\cdot12=144$ square inches, so the area is $\dfrac{18}{144}=\boxed{\dfrac{1}{8}}$ of a square foot.

The poster converted the legs to $\dfrac{3}{4}$ ft and $\dfrac{1}{3}$ ft correctly but multiplied them without the factor $\dfrac{1}{2}$: $\dfrac{3}{4}\cdot\dfrac{1}{3}=\dfrac{1}{4}$ is the area of the rectangle with these sides, twice the triangle's area $\dfrac{1}{8}$.
\begin{asy}
usepackage("amsmath");
size(4.5cm);
defaultpen(fontsize(10pt));
pair C=(0,0),A=(9,0),B=(0,4);
draw(A--B--C--cycle,linewidth(0.9));
draw((0,0.6)--(0.6,0.6)--(0.6,0));
label("$9$",(A+C)/2,S); label("$4$",(B+C)/2,W); label("$\sqrt{97}$",(A+B)/2,NE);
label("\mbox{}",(min(currentpicture,user=true)+max(currentpicture,user=true))/2,black); // keeps the text after the figure black when LaTeX typesets the labels inline
\end{asy}
